\documentclass{article}
\usepackage{iclr2027_conference,times}
\usepackage[T1]{fontenc}
%%%%% NEW MATH DEFINITIONS %%%%%

\usepackage{amsmath,amsfonts,bm}

% Mark sections of captions for referring to divisions of figures
\newcommand{\figleft}{{\em (Left)}}
\newcommand{\figcenter}{{\em (Center)}}
\newcommand{\figright}{{\em (Right)}}
\newcommand{\figtop}{{\em (Top)}}
\newcommand{\figbottom}{{\em (Bottom)}}
\newcommand{\captiona}{{\em (a)}}
\newcommand{\captionb}{{\em (b)}}
\newcommand{\captionc}{{\em (c)}}
\newcommand{\captiond}{{\em (d)}}

% Highlight a newly defined term
\newcommand{\newterm}[1]{{\bf #1}}


% Figure reference, lower-case.
\def\figref#1{figure~\ref{#1}}
% Figure reference, capital. For start of sentence
\def\Figref#1{Figure~\ref{#1}}
\def\twofigref#1#2{figures \ref{#1} and \ref{#2}}
\def\quadfigref#1#2#3#4{figures \ref{#1}, \ref{#2}, \ref{#3} and \ref{#4}}
% Section reference, lower-case.
\def\secref#1{section~\ref{#1}}
% Section reference, capital.
\def\Secref#1{Section~\ref{#1}}
% Reference to two sections.
\def\twosecrefs#1#2{sections \ref{#1} and \ref{#2}}
% Reference to three sections.
\def\secrefs#1#2#3{sections \ref{#1}, \ref{#2} and \ref{#3}}
% Reference to an equation, lower-case.
\def\eqref#1{equation~\ref{#1}}
% Reference to an equation, upper case
\def\Eqref#1{Equation~\ref{#1}}
% A raw reference to an equation---avoid using if possible
\def\plaineqref#1{\ref{#1}}
% Reference to a chapter, lower-case.
\def\chapref#1{chapter~\ref{#1}}
% Reference to an equation, upper case.
\def\Chapref#1{Chapter~\ref{#1}}
% Reference to a range of chapters
\def\rangechapref#1#2{chapters\ref{#1}--\ref{#2}}
% Reference to an algorithm, lower-case.
\def\algref#1{algorithm~\ref{#1}}
% Reference to an algorithm, upper case.
\def\Algref#1{Algorithm~\ref{#1}}
\def\twoalgref#1#2{algorithms \ref{#1} and \ref{#2}}
\def\Twoalgref#1#2{Algorithms \ref{#1} and \ref{#2}}
% Reference to a part, lower case
\def\partref#1{part~\ref{#1}}
% Reference to a part, upper case
\def\Partref#1{Part~\ref{#1}}
\def\twopartref#1#2{parts \ref{#1} and \ref{#2}}

\def\ceil#1{\lceil #1 \rceil}
\def\floor#1{\lfloor #1 \rfloor}
\def\1{\bm{1}}
\newcommand{\train}{\mathcal{D}}
\newcommand{\valid}{\mathcal{D_{\mathrm{valid}}}}
\newcommand{\test}{\mathcal{D_{\mathrm{test}}}}

\def\eps{{\epsilon}}


% Random variables
\def\reta{{\textnormal{$\eta$}}}
\def\ra{{\textnormal{a}}}
\def\rb{{\textnormal{b}}}
\def\rc{{\textnormal{c}}}
\def\rd{{\textnormal{d}}}
\def\re{{\textnormal{e}}}
\def\rf{{\textnormal{f}}}
\def\rg{{\textnormal{g}}}
\def\rh{{\textnormal{h}}}
\def\ri{{\textnormal{i}}}
\def\rj{{\textnormal{j}}}
\def\rk{{\textnormal{k}}}
\def\rl{{\textnormal{l}}}
% rm is already a command, just don't name any random variables m
\def\rn{{\textnormal{n}}}
\def\ro{{\textnormal{o}}}
\def\rp{{\textnormal{p}}}
\def\rq{{\textnormal{q}}}
\def\rr{{\textnormal{r}}}
\def\rs{{\textnormal{s}}}
\def\rt{{\textnormal{t}}}
\def\ru{{\textnormal{u}}}
\def\rv{{\textnormal{v}}}
\def\rw{{\textnormal{w}}}
\def\rx{{\textnormal{x}}}
\def\ry{{\textnormal{y}}}
\def\rz{{\textnormal{z}}}

% Random vectors
\def\rvepsilon{{\mathbf{\epsilon}}}
\def\rvtheta{{\mathbf{\theta}}}
\def\rva{{\mathbf{a}}}
\def\rvb{{\mathbf{b}}}
\def\rvc{{\mathbf{c}}}
\def\rvd{{\mathbf{d}}}
\def\rve{{\mathbf{e}}}
\def\rvf{{\mathbf{f}}}
\def\rvg{{\mathbf{g}}}
\def\rvh{{\mathbf{h}}}
\def\rvu{{\mathbf{i}}}
\def\rvj{{\mathbf{j}}}
\def\rvk{{\mathbf{k}}}
\def\rvl{{\mathbf{l}}}
\def\rvm{{\mathbf{m}}}
\def\rvn{{\mathbf{n}}}
\def\rvo{{\mathbf{o}}}
\def\rvp{{\mathbf{p}}}
\def\rvq{{\mathbf{q}}}
\def\rvr{{\mathbf{r}}}
\def\rvs{{\mathbf{s}}}
\def\rvt{{\mathbf{t}}}
\def\rvu{{\mathbf{u}}}
\def\rvv{{\mathbf{v}}}
\def\rvw{{\mathbf{w}}}
\def\rvx{{\mathbf{x}}}
\def\rvy{{\mathbf{y}}}
\def\rvz{{\mathbf{z}}}

% Elements of random vectors
\def\erva{{\textnormal{a}}}
\def\ervb{{\textnormal{b}}}
\def\ervc{{\textnormal{c}}}
\def\ervd{{\textnormal{d}}}
\def\erve{{\textnormal{e}}}
\def\ervf{{\textnormal{f}}}
\def\ervg{{\textnormal{g}}}
\def\ervh{{\textnormal{h}}}
\def\ervi{{\textnormal{i}}}
\def\ervj{{\textnormal{j}}}
\def\ervk{{\textnormal{k}}}
\def\ervl{{\textnormal{l}}}
\def\ervm{{\textnormal{m}}}
\def\ervn{{\textnormal{n}}}
\def\ervo{{\textnormal{o}}}
\def\ervp{{\textnormal{p}}}
\def\ervq{{\textnormal{q}}}
\def\ervr{{\textnormal{r}}}
\def\ervs{{\textnormal{s}}}
\def\ervt{{\textnormal{t}}}
\def\ervu{{\textnormal{u}}}
\def\ervv{{\textnormal{v}}}
\def\ervw{{\textnormal{w}}}
\def\ervx{{\textnormal{x}}}
\def\ervy{{\textnormal{y}}}
\def\ervz{{\textnormal{z}}}

% Random matrices
\def\rmA{{\mathbf{A}}}
\def\rmB{{\mathbf{B}}}
\def\rmC{{\mathbf{C}}}
\def\rmD{{\mathbf{D}}}
\def\rmE{{\mathbf{E}}}
\def\rmF{{\mathbf{F}}}
\def\rmG{{\mathbf{G}}}
\def\rmH{{\mathbf{H}}}
\def\rmI{{\mathbf{I}}}
\def\rmJ{{\mathbf{J}}}
\def\rmK{{\mathbf{K}}}
\def\rmL{{\mathbf{L}}}
\def\rmM{{\mathbf{M}}}
\def\rmN{{\mathbf{N}}}
\def\rmO{{\mathbf{O}}}
\def\rmP{{\mathbf{P}}}
\def\rmQ{{\mathbf{Q}}}
\def\rmR{{\mathbf{R}}}
\def\rmS{{\mathbf{S}}}
\def\rmT{{\mathbf{T}}}
\def\rmU{{\mathbf{U}}}
\def\rmV{{\mathbf{V}}}
\def\rmW{{\mathbf{W}}}
\def\rmX{{\mathbf{X}}}
\def\rmY{{\mathbf{Y}}}
\def\rmZ{{\mathbf{Z}}}

% Elements of random matrices
\def\ermA{{\textnormal{A}}}
\def\ermB{{\textnormal{B}}}
\def\ermC{{\textnormal{C}}}
\def\ermD{{\textnormal{D}}}
\def\ermE{{\textnormal{E}}}
\def\ermF{{\textnormal{F}}}
\def\ermG{{\textnormal{G}}}
\def\ermH{{\textnormal{H}}}
\def\ermI{{\textnormal{I}}}
\def\ermJ{{\textnormal{J}}}
\def\ermK{{\textnormal{K}}}
\def\ermL{{\textnormal{L}}}
\def\ermM{{\textnormal{M}}}
\def\ermN{{\textnormal{N}}}
\def\ermO{{\textnormal{O}}}
\def\ermP{{\textnormal{P}}}
\def\ermQ{{\textnormal{Q}}}
\def\ermR{{\textnormal{R}}}
\def\ermS{{\textnormal{S}}}
\def\ermT{{\textnormal{T}}}
\def\ermU{{\textnormal{U}}}
\def\ermV{{\textnormal{V}}}
\def\ermW{{\textnormal{W}}}
\def\ermX{{\textnormal{X}}}
\def\ermY{{\textnormal{Y}}}
\def\ermZ{{\textnormal{Z}}}

% Vectors
\def\vzero{{\bm{0}}}
\def\vone{{\bm{1}}}
\def\vmu{{\bm{\mu}}}
\def\vtheta{{\bm{\theta}}}
\def\va{{\bm{a}}}
\def\vb{{\bm{b}}}
\def\vc{{\bm{c}}}
\def\vd{{\bm{d}}}
\def\ve{{\bm{e}}}
\def\vf{{\bm{f}}}
\def\vg{{\bm{g}}}
\def\vh{{\bm{h}}}
\def\vi{{\bm{i}}}
\def\vj{{\bm{j}}}
\def\vk{{\bm{k}}}
\def\vl{{\bm{l}}}
\def\vm{{\bm{m}}}
\def\vn{{\bm{n}}}
\def\vo{{\bm{o}}}
\def\vp{{\bm{p}}}
\def\vq{{\bm{q}}}
\def\vr{{\bm{r}}}
\def\vs{{\bm{s}}}
\def\vt{{\bm{t}}}
\def\vu{{\bm{u}}}
\def\vv{{\bm{v}}}
\def\vw{{\bm{w}}}
\def\vx{{\bm{x}}}
\def\vy{{\bm{y}}}
\def\vz{{\bm{z}}}

% Elements of vectors
\def\evalpha{{\alpha}}
\def\evbeta{{\beta}}
\def\evepsilon{{\epsilon}}
\def\evlambda{{\lambda}}
\def\evomega{{\omega}}
\def\evmu{{\mu}}
\def\evpsi{{\psi}}
\def\evsigma{{\sigma}}
\def\evtheta{{\theta}}
\def\eva{{a}}
\def\evb{{b}}
\def\evc{{c}}
\def\evd{{d}}
\def\eve{{e}}
\def\evf{{f}}
\def\evg{{g}}
\def\evh{{h}}
\def\evi{{i}}
\def\evj{{j}}
\def\evk{{k}}
\def\evl{{l}}
\def\evm{{m}}
\def\evn{{n}}
\def\evo{{o}}
\def\evp{{p}}
\def\evq{{q}}
\def\evr{{r}}
\def\evs{{s}}
\def\evt{{t}}
\def\evu{{u}}
\def\evv{{v}}
\def\evw{{w}}
\def\evx{{x}}
\def\evy{{y}}
\def\evz{{z}}

% Matrix
\def\mA{{\bm{A}}}
\def\mB{{\bm{B}}}
\def\mC{{\bm{C}}}
\def\mD{{\bm{D}}}
\def\mE{{\bm{E}}}
\def\mF{{\bm{F}}}
\def\mG{{\bm{G}}}
\def\mH{{\bm{H}}}
\def\mI{{\bm{I}}}
\def\mJ{{\bm{J}}}
\def\mK{{\bm{K}}}
\def\mL{{\bm{L}}}
\def\mM{{\bm{M}}}
\def\mN{{\bm{N}}}
\def\mO{{\bm{O}}}
\def\mP{{\bm{P}}}
\def\mQ{{\bm{Q}}}
\def\mR{{\bm{R}}}
\def\mS{{\bm{S}}}
\def\mT{{\bm{T}}}
\def\mU{{\bm{U}}}
\def\mV{{\bm{V}}}
\def\mW{{\bm{W}}}
\def\mX{{\bm{X}}}
\def\mY{{\bm{Y}}}
\def\mZ{{\bm{Z}}}
\def\mBeta{{\bm{\beta}}}
\def\mPhi{{\bm{\Phi}}}
\def\mLambda{{\bm{\Lambda}}}
\def\mSigma{{\bm{\Sigma}}}

% Tensor
\DeclareMathAlphabet{\mathsfit}{\encodingdefault}{\sfdefault}{m}{sl}
\SetMathAlphabet{\mathsfit}{bold}{\encodingdefault}{\sfdefault}{bx}{n}
\newcommand{\tens}[1]{\bm{\mathsfit{#1}}}
\def\tA{{\tens{A}}}
\def\tB{{\tens{B}}}
\def\tC{{\tens{C}}}
\def\tD{{\tens{D}}}
\def\tE{{\tens{E}}}
\def\tF{{\tens{F}}}
\def\tG{{\tens{G}}}
\def\tH{{\tens{H}}}
\def\tI{{\tens{I}}}
\def\tJ{{\tens{J}}}
\def\tK{{\tens{K}}}
\def\tL{{\tens{L}}}
\def\tM{{\tens{M}}}
\def\tN{{\tens{N}}}
\def\tO{{\tens{O}}}
\def\tP{{\tens{P}}}
\def\tQ{{\tens{Q}}}
\def\tR{{\tens{R}}}
\def\tS{{\tens{S}}}
\def\tT{{\tens{T}}}
\def\tU{{\tens{U}}}
\def\tV{{\tens{V}}}
\def\tW{{\tens{W}}}
\def\tX{{\tens{X}}}
\def\tY{{\tens{Y}}}
\def\tZ{{\tens{Z}}}


% Graph
\def\gA{{\mathcal{A}}}
\def\gB{{\mathcal{B}}}
\def\gC{{\mathcal{C}}}
\def\gD{{\mathcal{D}}}
\def\gE{{\mathcal{E}}}
\def\gF{{\mathcal{F}}}
\def\gG{{\mathcal{G}}}
\def\gH{{\mathcal{H}}}
\def\gI{{\mathcal{I}}}
\def\gJ{{\mathcal{J}}}
\def\gK{{\mathcal{K}}}
\def\gL{{\mathcal{L}}}
\def\gM{{\mathcal{M}}}
\def\gN{{\mathcal{N}}}
\def\gO{{\mathcal{O}}}
\def\gP{{\mathcal{P}}}
\def\gQ{{\mathcal{Q}}}
\def\gR{{\mathcal{R}}}
\def\gS{{\mathcal{S}}}
\def\gT{{\mathcal{T}}}
\def\gU{{\mathcal{U}}}
\def\gV{{\mathcal{V}}}
\def\gW{{\mathcal{W}}}
\def\gX{{\mathcal{X}}}
\def\gY{{\mathcal{Y}}}
\def\gZ{{\mathcal{Z}}}

% Sets
\def\sA{{\mathbb{A}}}
\def\sB{{\mathbb{B}}}
\def\sC{{\mathbb{C}}}
\def\sD{{\mathbb{D}}}
% Don't use a set called E, because this would be the same as our symbol
% for expectation.
\def\sF{{\mathbb{F}}}
\def\sG{{\mathbb{G}}}
\def\sH{{\mathbb{H}}}
\def\sI{{\mathbb{I}}}
\def\sJ{{\mathbb{J}}}
\def\sK{{\mathbb{K}}}
\def\sL{{\mathbb{L}}}
\def\sM{{\mathbb{M}}}
\def\sN{{\mathbb{N}}}
\def\sO{{\mathbb{O}}}
\def\sP{{\mathbb{P}}}
\def\sQ{{\mathbb{Q}}}
\def\sR{{\mathbb{R}}}
\def\sS{{\mathbb{S}}}
\def\sT{{\mathbb{T}}}
\def\sU{{\mathbb{U}}}
\def\sV{{\mathbb{V}}}
\def\sW{{\mathbb{W}}}
\def\sX{{\mathbb{X}}}
\def\sY{{\mathbb{Y}}}
\def\sZ{{\mathbb{Z}}}

% Entries of a matrix
\def\emLambda{{\Lambda}}
\def\emA{{A}}
\def\emB{{B}}
\def\emC{{C}}
\def\emD{{D}}
\def\emE{{E}}
\def\emF{{F}}
\def\emG{{G}}
\def\emH{{H}}
\def\emI{{I}}
\def\emJ{{J}}
\def\emK{{K}}
\def\emL{{L}}
\def\emM{{M}}
\def\emN{{N}}
\def\emO{{O}}
\def\emP{{P}}
\def\emQ{{Q}}
\def\emR{{R}}
\def\emS{{S}}
\def\emT{{T}}
\def\emU{{U}}
\def\emV{{V}}
\def\emW{{W}}
\def\emX{{X}}
\def\emY{{Y}}
\def\emZ{{Z}}
\def\emSigma{{\Sigma}}

% entries of a tensor
% Same font as tensor, without \bm wrapper
\newcommand{\etens}[1]{\mathsfit{#1}}
\def\etLambda{{\etens{\Lambda}}}
\def\etA{{\etens{A}}}
\def\etB{{\etens{B}}}
\def\etC{{\etens{C}}}
\def\etD{{\etens{D}}}
\def\etE{{\etens{E}}}
\def\etF{{\etens{F}}}
\def\etG{{\etens{G}}}
\def\etH{{\etens{H}}}
\def\etI{{\etens{I}}}
\def\etJ{{\etens{J}}}
\def\etK{{\etens{K}}}
\def\etL{{\etens{L}}}
\def\etM{{\etens{M}}}
\def\etN{{\etens{N}}}
\def\etO{{\etens{O}}}
\def\etP{{\etens{P}}}
\def\etQ{{\etens{Q}}}
\def\etR{{\etens{R}}}
\def\etS{{\etens{S}}}
\def\etT{{\etens{T}}}
\def\etU{{\etens{U}}}
\def\etV{{\etens{V}}}
\def\etW{{\etens{W}}}
\def\etX{{\etens{X}}}
\def\etY{{\etens{Y}}}
\def\etZ{{\etens{Z}}}

% The true underlying data generating distribution
\newcommand{\pdata}{p_{\rm{data}}}
% The empirical distribution defined by the training set
\newcommand{\ptrain}{\hat{p}_{\rm{data}}}
\newcommand{\Ptrain}{\hat{P}_{\rm{data}}}
% The model distribution
\newcommand{\pmodel}{p_{\rm{model}}}
\newcommand{\Pmodel}{P_{\rm{model}}}
\newcommand{\ptildemodel}{\tilde{p}_{\rm{model}}}
% Stochastic autoencoder distributions
\newcommand{\pencode}{p_{\rm{encoder}}}
\newcommand{\pdecode}{p_{\rm{decoder}}}
\newcommand{\precons}{p_{\rm{reconstruct}}}

\newcommand{\laplace}{\mathrm{Laplace}} % Laplace distribution

\newcommand{\E}{\mathbb{E}}
\newcommand{\Ls}{\mathcal{L}}
\newcommand{\R}{\mathbb{R}}
\newcommand{\emp}{\tilde{p}}
\newcommand{\lr}{\alpha}
\newcommand{\reg}{\lambda}
\newcommand{\rect}{\mathrm{rectifier}}
\newcommand{\softmax}{\mathrm{softmax}}
\newcommand{\sigmoid}{\sigma}
\newcommand{\softplus}{\zeta}
\newcommand{\KL}{D_{\mathrm{KL}}}
\newcommand{\Var}{\mathrm{Var}}
\newcommand{\standarderror}{\mathrm{SE}}
\newcommand{\Cov}{\mathrm{Cov}}
% Wolfram Mathworld says $L^2$ is for function spaces and $\ell^2$ is for vectors
% But then they seem to use $L^2$ for vectors throughout the site, and so does
% wikipedia.
\newcommand{\normlzero}{L^0}
\newcommand{\normlone}{L^1}
\newcommand{\normltwo}{L^2}
\newcommand{\normlp}{L^p}
\newcommand{\normmax}{L^\infty}

\newcommand{\parents}{Pa} % See usage in notation.tex. Chosen to match Daphne's book.

\DeclareMathOperator*{\argmax}{arg\,max}
\DeclareMathOperator*{\argmin}{arg\,min}

\DeclareMathOperator{\sign}{sign}
\DeclareMathOperator{\Tr}{Tr}
\let\ab\allowbreak

\usepackage{hyperref}
\usepackage{url}
\usepackage{amsmath,amssymb,amsthm}
\usepackage{booktabs}
\usepackage{graphicx}
\usepackage{xcolor}
\usepackage{enumitem}
\usepackage{multirow}

% ---------------------------------------------------------------- numbers
% Every number in the text is a macro. numbers.tex (generated from the result files) is read first and wins;
% numbers_prior.tex (generated by exp/prior_stats.py) fills derived statistics of the prior grid;
% numbers_placeholder.tex (\providecommand only) fills whatever is still undefined, with \TBD for open cells.
\InputIfFileExists{numbers.tex}{}{}
\InputIfFileExists{numbers_prior.tex}{}{}
% numbers_placeholder.tex -- fallback values for every number in main.tex.
%
% Load order in main.tex:  numbers.tex (if it exists; generated from result files, use \newcommand or \def)
%                          -> numbers_prior.tex (if it exists; generated by exp/prior_stats.py)
%                          -> this file.
% Everything here is \providecommand, so a macro defined by an earlier file wins and this file only fills
% gaps. Values below are the prior-grid values (results/THINKING_VS_DATA.md, results/GATE_FINETUNE_STUDY.md,
% results/LOGIC_FINETUNE.md) or \TBD for cells preregistered in PREDICTIONS.md and not yet run.
% Percentages are written without the % sign; the text adds \%.

\providecommand{\TBD}{\textbf{[TBD]}}

% ---------------------------------------------------------------- setup
\providecommand{\nTrain}{180}
\providecommand{\nTrainSmall}{60}
\providecommand{\nTrainLarge}{540}
\providecommand{\nTest}{240}
\providecommand{\nTestClass}{120}
\providecommand{\loraRank}{8}
\providecommand{\loraAlpha}{16}
\providecommand{\loraDropout}{0.05}
\providecommand{\nEpochs}{3}
\providecommand{\learningRate}{2\times 10^{-4}}
\providecommand{\ciHalf}{6}
\providecommand{\nShots}{4}
\providecommand{\nForms}{12}
\providecommand{\nFormsHalf}{6}
\providecommand{\maxNewPrime}{512}
\providecommand{\maxNewDiv}{176}
\providecommand{\maxNewValid}{320}
\providecommand{\nPriorRuns}{36}
\providecommand{\priorWallMin}{319}
\providecommand{\priorVerdictDate}{2026-08-04}
\providecommand{\alphaLevel}{0.05}

% ---------------------------------------------------------------- trace structure and training tokens
\providecommand{\divsevenSteps}{4}
\providecommand{\divsevenSixSteps}{6}
\providecommand{\divsevenTableSize}{70}
\providecommand{\divsevenStepSamples}{720}
\providecommand{\primeStepsMin}{11}
\providecommand{\primeStepsMax}{25}
\providecommand{\primeStepsMeanPrimes}{19.2}
\providecommand{\primeStepsMedComp}{1}
\providecommand{\primeStepsMeanComp}{2.6}
\providecommand{\primeTraceTok}{71.5}
\providecommand{\divsevenTraceTok}{58.0}
\providecommand{\validTraceTok}{37.0}
\providecommand{\tokA}{360}
\providecommand{\primeTokB}{13{,}300}
\providecommand{\divsevenTokB}{10{,}800}
\providecommand{\validTokB}{7{,}000}
\providecommand{\tokRatioLow}{19}
\providecommand{\tokRatioHigh}{37}
\providecommand{\runMinLow}{2}
\providecommand{\runMinHigh}{17}
\providecommand{\validPool}{1{,}600}
\providecommand{\cTokDiff}{6}
\providecommand{\baseEvalRatio}{38}
\providecommand{\nProbeFeatures}{40}

% ---------------------------------------------------------------- constants of the tasks (exact, not results)
\providecommand{\condLow}{0.10}
\providecommand{\condHigh}{0.22}
\providecommand{\condMarg}{0.14}
\providecommand{\hardNegPop}{996}
\providecommand{\compPop}{7{,}939}
\providecommand{\hardNegFrac}{12.5}
\providecommand{\winMargin}{10}
\providecommand{\neitherMargin}{8}

% ---------------------------------------------------------------- surface probe and rules (prime test set)
% surface probes: experiments/analysis_probe.py (results/analysis/probe.json). prime values are its seed-0
% output (best by 5-fold CV = logreg on one-hot + engineered features; one-hot logreg alone). div7 and valid
% are stand-ins from a plain one-hot / atom-renamed n-gram logreg (C = 1) until analysis_probe.py covers them.
\providecommand{\primeProbe}{89.2}
\providecommand{\primeProbeOneHot}{78.8}
\providecommand{\divsevenProbe}{50.0}
\providecommand{\validProbe}{100.0}
% rules evaluated on this paper's prime test set (where A = 82.3)
\providecommand{\ruleOdd}{76.2}
\providecommand{\ruleNoTwoThree}{86.7}
\providecommand{\ruleNoSeven}{93.8}
% the same rules on the earlier answer-only study's prime test set (where answers gave 82.9)
\providecommand{\ruleOddPrior}{78.8}
\providecommand{\ruleNoTwoThreePrior}{89.2}
\providecommand{\primeAagreeOdd}{93.5}
\providecommand{\ruleLastDigit}{82.5}
\providecommand{\primeAagreeLastDigitLo}{99.6}
\providecommand{\primeAagreeLastDigitHi}{100.0}
\providecommand{\primeAprimesCorrect}{99.6}
\providecommand{\primeAcompCorrect}{65.0}
\providecommand{\primeHardNegN}{15}
\providecommand{\primeHardNegAcorrect}{0 and 0}
\providecommand{\primeBprimesCorrect}{19.6}
\providecommand{\primeBcompCorrect}{95.8}
\providecommand{\primeBlargePrimes}{46.7}
\providecommand{\primeBlargeComp}{90.8}

% ---------------------------------------------------------------- prime (Qwen2.5-1.5B, n = 180 unless noted)
\providecommand{\primeBase}{59.6}
\providecommand{\primeA}{82.3}
\providecommand{\primeAsZero}{82.1}
\providecommand{\primeAsOne}{82.5}
\providecommand{\primeB}{57.7}
\providecommand{\primeBsZero}{57.9}
\providecommand{\primeBsOne}{57.5}
\providecommand{\primeC}{52.1}
\providecommand{\primeCsZero}{50.0}
\providecommand{\primeCsOne}{54.2}
\providecommand{\primeAsmall}{59.6}
\providecommand{\primeAlarge}{75.0}
\providecommand{\primeBsmall}{57.1}
\providecommand{\primeBlarge}{68.8}
\providecommand{\primeAtok}{80.8}
\providecommand{\primeAtokN}{1{,}320}
\providecommand{\primeBorn}{61.0}
\providecommand{\primeBornN}{9}
\providecommand{\primeBApsZero}{6\times 10^{-7}}
\providecommand{\primeBApsOne}{2\times 10^{-7}}
\providecommand{\primeBCpsZero}{0.22}
\providecommand{\primeBCpsOne}{0.33}

% ---------------------------------------------------------------- div7 (4-digit)
\providecommand{\divsevenBase}{62.5}
\providecommand{\divsevenBaseYes}{84.2}
\providecommand{\divsevenBaseNo}{40.8}
\providecommand{\divsevenA}{51.5}
\providecommand{\divsevenAsZero}{52.1}
\providecommand{\divsevenAsOne}{50.8}
\providecommand{\divsevenAyessZero}{25.8}
\providecommand{\divsevenAnosZero}{78.3}
\providecommand{\divsevenAyessOne}{82.5}
\providecommand{\divsevenAnosOne}{19.2}
\providecommand{\divsevenB}{92.5}
\providecommand{\divsevenBsZero}{92.5}
\providecommand{\divsevenByes}{89.2}
\providecommand{\divsevenBno}{95.8}
\providecommand{\divsevenC}{50.0}
\providecommand{\divsevenCsZero}{49.6}
\providecommand{\divsevenCsOne}{50.4}
\providecommand{\divsevenAsmall}{50.0}
\providecommand{\divsevenAlarge}{50.0}
\providecommand{\divsevenBsmall}{50.0}
\providecommand{\divsevenBApsZero}{4\times 10^{-27}}
\providecommand{\divsevenBApsZeroDisc}{99/2}
\providecommand{\divsevenBCpsZero}{1\times 10^{-25}}
\providecommand{\divsevenABasepsZero}{0.047}
\providecommand{\divsevenABasepsOne}{0.0046}

% ---------------------------------------------------------------- valid (propositional validity, Z3 labels)
\providecommand{\validBase}{30.0}
\providecommand{\validA}{99.6}
\providecommand{\validAsZero}{99.2}
\providecommand{\validAsOne}{100.0}
\providecommand{\validB}{90.0}
\providecommand{\validBsZero}{100.0}
\providecommand{\validBsOne}{80.0}
\providecommand{\validC}{52.5}
\providecommand{\validCsZero}{54.2}
\providecommand{\validCsOne}{50.8}
\providecommand{\validAsmall}{97.5}
\providecommand{\validBsmall}{87.5}
\providecommand{\validAlarge}{100.0}
\providecommand{\validBlarge}{100.0}
\providecommand{\validAtok}{100.0}
\providecommand{\validBApsZero}{0.50}
\providecommand{\validBApsOne}{7\times 10^{-15}}
\providecommand{\validBCpsZero}{2\times 10^{-33}}
\providecommand{\validBCpsOne}{5\times 10^{-10}}

% ---------------------------------------------------------------- earlier answer-only study (shuffled-label control)
\providecommand{\priorTrainN}{240}
\providecommand{\shufMin}{48.5}
\providecommand{\shufMax}{52.5}
\providecommand{\priorCorrectMax}{99.8}
\providecommand{\priorDivthreeBase}{55.0}
\providecommand{\priorDivthreeA}{67.3}
\providecommand{\priorDivthreeShuf}{49.6}
\providecommand{\priorDivsevenBase}{53.3}
\providecommand{\priorDivsevenA}{51.0}
\providecommand{\priorDivsevenShuf}{52.5}
\providecommand{\priorDivthirteenBase}{54.2}
\providecommand{\priorDivthirteenA}{50.0}
\providecommand{\priorDivthirteenShuf}{50.0}
\providecommand{\priorSquareBase}{51.2}
\providecommand{\priorSquareA}{83.3}
\providecommand{\priorSquareShuf}{51.5}
\providecommand{\priorPrimeBase}{50.0}
\providecommand{\priorPrimeA}{82.9}
\providecommand{\priorPrimeShuf}{50.0}
\providecommand{\priorValidBase}{66.7}
\providecommand{\priorValidA}{99.4}
\providecommand{\priorValidShuf}{50.0}
\providecommand{\priorConsistBase}{50.0}
\providecommand{\priorConsistA}{99.8}
\providecommand{\priorConsistShuf}{48.5}

% ================================================================ preregistered cells (PREDICTIONS.md), not yet run
% S1: div7 replication
\providecommand{\divsevenBsOne}{\TBD}
\providecommand{\divsevenBsTwo}{\TBD}
\providecommand{\divsevenAsTwo}{\TBD}
\providecommand{\divsevenBApsOne}{\TBD}
\providecommand{\divsevenBApsTwo}{\TBD}
\providecommand{\verdictSOne}{\TBD}
% S2/S3: divisor family (4-digit, k = 4); mean of two seeds
\providecommand{\divtwoA}{\TBD}
\providecommand{\divtwoB}{\TBD}
\providecommand{\divtwoProbe}{\TBD}
\providecommand{\divthreeA}{\TBD}
\providecommand{\divthreeB}{\TBD}
\providecommand{\divthreeProbe}{\TBD}
\providecommand{\divelevenA}{\TBD}
\providecommand{\divelevenB}{\TBD}
\providecommand{\divelevenProbe}{\TBD}
\providecommand{\divthirteenA}{\TBD}
\providecommand{\divthirteenB}{\TBD}
\providecommand{\divthirteenProbe}{\TBD}
\providecommand{\verdictSTwo}{\TBD}
\providecommand{\verdictSThree}{\TBD}
% S4: scale
\providecommand{\threeBdivsevenA}{\TBD}
\providecommand{\threeBdivsevenB}{\TBD}
\providecommand{\threeBprimeA}{\TBD}
\providecommand{\threeBprimeB}{\TBD}
\providecommand{\threeBvalidA}{\TBD}
\providecommand{\threeBvalidB}{\TBD}
\providecommand{\halfBdivsevenA}{\TBD}
\providecommand{\halfBdivsevenB}{\TBD}
\providecommand{\sevenBdivsevenA}{\TBD}
\providecommand{\sevenBdivsevenB}{\TBD}
\providecommand{\verdictSFour}{\TBD}
% S5: length test (6-digit div7, k = 6)
\providecommand{\stepAcc}{\TBD}
\providecommand{\divsevenSixPred}{\TBD}
\providecommand{\divsevenSixPredBal}{\TBD}
\providecommand{\divsevenSixA}{\TBD}
\providecommand{\divsevenSixB}{\TBD}
\providecommand{\divsevenSixByes}{\TBD}
\providecommand{\verdictSFive}{\TBD}
% S6: prime hard negatives (odd composites with no prime factor <= 7)
\providecommand{\primeHardN}{\TBD}
\providecommand{\primeHardA}{\TBD}
\providecommand{\primeHardB}{\TBD}
\providecommand{\verdictSSix}{\TBD}
% answer-then-trace (B') and test-time-compute controls on the base model
\providecommand{\divsevenBp}{\TBD}
\providecommand{\primeCot}{\TBD}
\providecommand{\divsevenCot}{\TBD}
\providecommand{\validCot}{\TBD}
\providecommand{\primeFew}{\TBD}
\providecommand{\divsevenFew}{\TBD}
\providecommand{\validFew}{\TBD}


\theoremstyle{definition}
\newtheorem{remark}{Remark}
\newcommand{\pp}{\,pp}
% figure if present (figures/ or paper root), otherwise a framed placeholder of the given height
\newcommand{\figorbox}[3]{%
  \IfFileExists{figures/#1}{\includegraphics[width=#2]{figures/#1}}{%
  \IfFileExists{#1}{\includegraphics[width=#2]{#1}}{%
  \fbox{\parbox[c][#3][c]{#2}{\centering\small\textit{Figure pending}}}}}}

\title{Perfect Labels for Free? Exact Labels Teach\\ Shortcuts, Exact Traces Teach Short Programs}

% Submission is double blind; the author block is filled in for the camera-ready only.
\author{Anonymous authors\\ Paper under double-blind review}

%\iclrfinalcopy % camera-ready only
\begin{document}
\maketitle

\begin{abstract}
Exact verifiers (primality tests, SAT solvers, unit tests) label data with no noise at no marginal cost, and they supply the rewards of verifiable-reward post-training. We use them as a measurement instrument: with label noise removed, the effect of the supervision \emph{format} can be isolated. We fine-tune Qwen2.5-1.5B-Instruct with LoRA on \nTrain{} exactly labeled examples per task and compare supervision by the answer alone with supervision by a deterministic procedure trace followed by the answer, against a length-matched scrambled-trace control. Answer-only training recovers about what simple probes on the input's surface recover. On 4-digit primality it reaches \primeA\%, between a one-hot linear probe (\primeProbeOneHot\%) and the best cross-validated probe (\primeProbe\%); its answers coincide with the rule ``the last digit is 1, 3, 7 or 9'' on \primeAagreeLastDigitLo--\primeAagreeLastDigitHi\% of test items, and it misclassifies all \primeHardNegN{} odd composites without a factor $\le 7$ in the test set. On 4-digit divisibility by 7, which has no surface rule, answers stay at chance (\divsevenA\%) while a \divsevenSteps-step long-division trace reaches \divsevenB\% and its scrambled counterpart \divsevenC\%. On primality, whose trial-division trace runs up to \primeStepsMax{} steps, traces hurt (\primeB\%). The results fit a simple map: answer supervision works when the task is surface-learnable, trace supervision works when the per-step accuracy compounded over the trace length stays high. We preregister tests of the map (a divisor family, a 6-digit length test, hard negatives, an answer-then-trace arm, prompting controls, a larger model) and report every cell. The lesson concerns how to design supervision for verifiable-reward pipelines, not how to do arithmetic.
\end{abstract}

\section{Introduction}
Verifiable-reward post-training rests on programs that decide correctness exactly: answer checkers for mathematics, unit tests for code, solvers for logic \citep{lambert2024tulu,deepseek2025r1}. Such verifiers are often described as a source of perfect labels for free. They are perfect in one sense: a verified label is never wrong. Whether a model trained on those labels learns the computation the verifier performs, or a cheaper function that agrees with it on the training distribution, is a separate question, and in practice it is confounded with label noise, with task difficulty, and with the format of the supervision.

We use exact verifiers to remove the first confound. When every label is correct, any difference between two training regimes on the same inputs is a difference in what the supervision exposes. We choose small decision tasks whose structure is known in advance (which surface rules approximate the label, how many steps the natural procedure takes) so that what should be learnable can be stated before training. The point is not to teach a language model what a ten-line program already computes. It is that on such tasks we can tell which function the model learned, which is what a verifiable-reward pipeline for mathematics or code needs to know about its own supervision and usually cannot measure.

We compare two supervision formats on Qwen2.5-1.5B-Instruct. \emph{Answer-only} supervision (arm A) teaches a shortcut. On 4-digit primality the answer-trained model reaches \primeA\%, between a one-hot linear probe on the digits (\primeProbeOneHot\%) and the best cross-validated surface probe (\primeProbe\%). Its answers are a one-token rule: they coincide with ``the last digit is 1, 3, 7 or 9'' on \primeAagreeLastDigitLo--\primeAagreeLastDigitHi\% of test items, and that rule alone scores \ruleLastDigit\%. On divisibility by 7, where no surface rule exists, it stays at chance. \emph{Trace} supervision (arm B: a deterministic procedure trace, then the answer) succeeds when the trace is short. The \divsevenSteps-step long-division trace lifts divisibility by 7 from \divsevenA\% to \divsevenB\% (McNemar $p=\divsevenBApsZero$), while a scrambled trace with the same tokens (arm C) stays at \divsevenC\%. On primality, where the trial-division trace of a prime has \primeStepsMin--\primeStepsMax{} steps, traces hurt (\primeB\% vs.\ \primeA\%), and the loss sits where the trace is long: \primeBprimesCorrect\% of primes are classified correctly against \primeBcompCorrect\% of composites, whose trace usually stops after \primeStepsMedComp{} step.

These observations fit a two-condition map (Figure~\ref{fig:map}). Answer supervision works if the label is \emph{surface-learnable}, which a probe measures before any fine-tuning. Trace supervision works if the trace is \emph{executable}: with per-step accuracy $p$ and $k$ steps, the chance of an error-free trace, about $p^k$, must stay high. The map makes predictions we had not tested when we wrote it down, and we preregistered them (\S\ref{sec:prereg}): varying the divisor at fixed $k$, lengthening the input to $k=\divsevenSixSteps$, hard negatives for the primality shortcut, an answer-then-trace arm that separates the trace as training signal from the trace as test-time computation, prompting controls on the base model, and a larger model.

\begin{figure}[t]
\centering
\figorbox{fig_map.pdf}{0.62\linewidth}{5.2cm}
\caption{The map. Each task is placed by the accuracy of the surface probe trained on answer labels (horizontal) and by the predicted trace accuracy $p^k$ (vertical), where $p$ is the per-step accuracy of arm-B generations and $k$ the number of steps; color marks the winning arm (A, B, both, neither). The map predicts A to win on the right, B at the top, both at top right and neither at bottom left.}
\label{fig:map}
\end{figure}

\paragraph{Contributions.}
\begin{itemize}[leftmargin=*,itemsep=1pt,topsep=2pt]
\item A controlled comparison of supervision formats under exact labels (answers, procedure traces, length-matched scrambled traces, shuffled labels) with paired tests on problem-disjoint test sets (\S\ref{sec:setup}).
\item Evidence that exact answer labels teach the surface probe's function rather than the verifier's: probe-level accuracy, rule agreement and failure on hard negatives (\S\ref{sec:shortcut}).
\item Evidence that exact traces teach short programs: the trace unlocks divisibility by 7 through its content, not its tokens, and fails on primality with the per-class signature of errors compounding over trace length (\S\ref{sec:traces}).
\item A predictive map with preregistered tests whose results are all reported (Table~\ref{tab:prereg}), and a remark relating it to automata and to known results on scratchpads (\S\ref{sec:why}).
\end{itemize}
One result reversed an earlier conclusion of ours. A verdict written on \priorVerdictDate{} from a \nTrainSmall-example trace run stated that traces do not unlock divisibility by 7; the preregistered \nTrain-example cell, which finished afterwards with the same code, gave \divsevenB\%. We report both, and the replication of that cell was made a go/no-go condition for this submission before any new run started (\S\ref{sec:prereg}).

\section{Related work}
\paragraph{Intermediate steps as supervision.} Scratchpads \citep{nye2021show} and chain-of-thought (CoT) prompting \citep{wei2022chain} showed that emitting intermediate steps improves multi-step computation, and \citet{sprague2025cot} find the benefit of CoT concentrated on mathematical and symbolic tasks. For small transformers trained from scratch, \citet{lee2024teaching} show that the format of arithmetic data, including step-by-step formats, changes sample efficiency sharply, and \citet{charton2024gcd} explains a transformer's GCD predictions by rules built from divisors of the number base, the kind of surface rule we probe for. \citet{dziri2023faith} show how errors compound across the steps of compositional tasks. We fine-tune a pretrained model, vary only the completion, and hold labels exact.

\paragraph{Theory.} Parity-like targets give gradient-based learners almost no signal from end labels \citep{shalevshwartz2017failures}. Sub-task decomposition makes otherwise unlearnable sequence-to-sequence tasks learnable \citep{wies2023subtask}; CoT raises the expressivity of bounded-depth transformers \citep{feng2023towards,li2024chain}; \citet{abbe2024far} define the globality degree of a target, show that it bounds what transformers learn efficiently, and show that scratchpads, in particular an inductive scratchpad, can break that barrier; \citet{kim2025parity} prove that a transformer trained with teacher-forced intermediate parities learns parity efficiently. Section~\ref{sec:why} reads divisibility through this lens without extending these results.

\paragraph{What in a trace matters.} \citet{valmeekam2026beyond} train transformers from scratch on formally verifiable traces and find the semantics of the trace less important than usually assumed: correct-trace training yields invalid traces with correct answers, and unrelated traces can still help. Our scrambled control is a coarser intervention (token multiset kept, order destroyed); on divisibility it removes the entire gain. Process supervision rewards steps rather than outcomes \citep{lightman2024verify}, and program-generated logic corpora improve LLM reasoning \citep{morishita2024logic}.

\paragraph{Verifiers in post-training.} STaR fine-tunes on self-generated rationales that reach verified answers \citep{zelikman2022star}; reinforcement learning with verifiable rewards (RLVR) drives recent reasoning models \citep{lambert2024tulu,deepseek2025r1}. Fine-tuning on model-generated correct solutions can amplify spurious correlations, which per-step credit from negatives mitigates \citep{setlur2024rl}; SFT and RL generalize differently \citep{chu2025sft}; and RLVR mostly reweights what the base model already samples \citep{yue2025rl}. Our answer arm is an instance of shortcut learning \citep{geirhos2020shortcut}; exact labels do not prevent it.

\section{Setup}
\label{sec:setup}
\paragraph{Exact verifiers as an instrument.} Each task has an exact decision procedure: trial division for primality, the remainder for divisibility, and Z3 \citep{demoura2008z3} for validity (premises entail $C$ iff premises $\wedge\,\neg C$ is unsatisfiable). Labels are correct by construction, every arm sees the same prompts and labels, and arms differ only in the completion the model is trained to produce. The tasks are chosen so that the cheap functions that approximate each label are known.

\paragraph{Tasks.} Table~\ref{tab:tasks} summarizes them. \textbf{prime}: is a uniformly drawn 4-digit $n$ prime? The trace performs trial division by the primes up to $\lfloor\sqrt n\rfloor$, one ``$n \bmod p = m$.'' step per prime, and stops at the first divisor, so a prime needs \primeStepsMin--\primeStepsMax{} steps (mean \primeStepsMeanPrimes) and a composite a median of \primeStepsMedComp. \textbf{div7}: is a 4-digit $n$ divisible by 7? The trace performs long division, $r\leftarrow(10r+x_t)\bmod 7$ for each digit $x_t$, so $k=\divsevenSteps$ for every item. \textbf{valid}: does a propositional conclusion follow from its premises? Arguments instantiate \nForms{} forms (\nFormsHalf{} valid, \nFormsHalf{} fallacies) with random atom letters; the trace is the Z3 certificate, a countermodel or an unsatisfiable core. Train and test share forms but not instances, so this task is surface-learnable by design and serves as a check (Appendix~\ref{app:logic}). Each test set has \nTest{} items (\nTestClass{} per class), is problem-disjoint from training, and is identical across arms; training sets are balanced samples of $n\in\{\nTrainSmall,\nTrain,\nTrainLarge\}$ from a training pool.

\begin{table}[t]
\caption{Tasks. $k$: steps in the trace (prime: for primes; a composite's trace stops at its smallest prime factor, after a median of \primeStepsMedComp{} step). Tokens: mean whitespace tokens per trace. Probe: the surface probe selected by 5-fold cross-validation on the arm-A training items (\S\ref{sec:shortcut}), test accuracy (\%) at $n=\nTrain$.}
\label{tab:tasks}
\centering\small
\begin{tabular}{llccc}
\toprule
task & trace procedure & $k$ & tokens & probe \\
\midrule
prime & trial division by the primes $\le\lfloor\sqrt n\rfloor$ & \primeStepsMin--\primeStepsMax & \primeTraceTok & \primeProbe \\
div7 & long division, $r\leftarrow(10r+x_t)\bmod 7$ & \divsevenSteps & \divsevenTraceTok & \divsevenProbe \\
valid & Z3 countermodel or unsatisfiable core & certificate & \validTraceTok & \validProbe \\
\bottomrule
\end{tabular}
\end{table}

\paragraph{Arms.} \textbf{base}: the model zero-shot. \textbf{A} (answers): the completion is ``Answer: Yes'' or ``Answer: No''. \textbf{B} (trace, then answer): the deterministic trace, a newline, then the answer (examples in Appendix~\ref{app:formats}). \textbf{C} (scrambled trace): B's traces are first permuted across training examples, so that trace length cannot signal the label (a composite's or an invalid argument's trace is shorter), then token-shuffled within each trace; the answer line keeps the example's own true label. Over the training pool, C has exactly B's multiset of whitespace tokens; in the sampled training sets the totals differ by at most \cTokDiff\%. \textbf{A-full}: answers only, on the whole training pool (\primeAtokN{} examples for prime, \validPool{} for valid). \textbf{Shuffled labels}: from an earlier study with the same model and LoRA configuration (\priorTrainN{} training items, answer-only prompts), with labels permuted across training items. Preregistered additions: \textbf{B$'$} (answer, then trace) on div7, in which the answer is emitted before the trace and so cannot be computed by it; and, on the base model only, a zero-shot CoT prompt and \nShots-shot trace demonstrations.

\paragraph{Training and evaluation.} LoRA \citep{hu2022lora} with $r=\loraRank$, $\alpha=\loraAlpha$, dropout \loraDropout{} on the query, key, value and output projections of Qwen2.5-1.5B-Instruct \citep{qwen2024qwen25}; AdamW, learning rate $\learningRate$, batch size 1, \nEpochs{} epochs, bf16; loss on completion tokens only, prompts in the chat template. Evaluation is greedy generation (at most \maxNewPrime, \maxNewDiv{} and \maxNewValid{} new tokens for prime, div7 and valid), and one extractor serves every arm: the last ``Answer: Yes/No'', else the first standalone yes/no. Each run is a fresh process on one 8\,GB consumer GPU and takes \runMinLow--\runMinHigh{} minutes at $n=\nTrain$.

\paragraph{Statistics, preregistration and provenance.}
\label{sec:prereg}
Wilson 95\% intervals at \nTest{} test items are about $\pm\ciHalf$\pp. Paired comparisons use the exact McNemar test \citep{mcnemar1947note} on per-item correctness within a seed, with Holm correction \citep{holm1979simple} across tasks within each hypothesis family at $\alpha=\alphaLevel$. Two preregistrations govern the paper. The first was frozen before the prior grid of \nPriorRuns{} runs and fixed arms, tasks and four endpoints: B beats A at matched examples (H1), B beats C (H2), B is more efficient per training token (H3), and on div7 B reaches about 60\% with A at chance (H4). The second (Appendix~\ref{app:prereg}) was committed before any run it lists and fixes the decision rules S1--S6 of Table~\ref{tab:prereg}; a ``win'' is a difference of at least \winMargin\pp{} with both seeds agreeing in sign. The surface probe belongs to neither: we defined it after seeing the prior grid and apply the same definition to every new task. \emph{Provenance.} The prior study's written verdict (\priorVerdictDate) stated that the trace does not unlock div7. It was written from the \nTrainSmall-example trace run alone (\divsevenBsmall\%) while the preregistered \nTrain-example cell was still running; that cell then finished at \divsevenB\% with no code change. The second preregistration made its replication (S1) a go/no-go condition for this submission.

\section{Results}
\label{sec:results}

\begin{table}[t]
\caption{Accuracy (\%) at $n=\nTrain$, Qwen2.5-1.5B-Instruct, \nTest{} test items per task (Wilson half-width about \ciHalf\pp). A: answers; B: trace, then answer; C: scrambled trace. A, B, C: mean of seeds 0 and 1, except div7 B (seed 0; seeds 1--2 are preregistered test S1). $p$: exact McNemar, per seed. Per-seed values in Appendix~\ref{app:tables}.}
\label{tab:main}
\centering\small
\begin{tabular}{lcccccll}
\toprule
task & base & probe & A & B & C & B vs.\ A, $p$ & B vs.\ C, $p$ \\
\midrule
prime & \primeBase & \primeProbe & \textbf{\primeA} & \primeB & \primeC & $\primeBApsZero$, $\primeBApsOne$ & $\primeBCpsZero$, $\primeBCpsOne$ \\
div7 & \divsevenBase & \divsevenProbe & \divsevenA & \textbf{\divsevenB} & \divsevenC & $\divsevenBApsZero$ & $\divsevenBCpsZero$ \\
valid & \validBase & \validProbe & \textbf{\validA} & \validB & \validC & $\validBApsZero$, $\validBApsOne$ & $\validBCpsZero$, $\validBCpsOne$ \\
\bottomrule
\end{tabular}
\end{table}



\subsection{Exact answer labels teach the surface}
\label{sec:shortcut}
\paragraph{Answers track surface probes.} We fit simple probes on the same \nTrain{} training items as arm A and score them on the test set: logistic regression, a one-hidden-layer MLP and $k$-nearest neighbours, on \nProbeFeatures{} one-hot digit-by-position features alone or together with engineered divisibility features (last digit even; last digit 0 or 5; digit sum mod 3 and mod 9; alternating digit sum mod 11). For arguments the features are word and character $n$-grams, with atoms renamed in order of first appearance. Hyperparameters and the reported probe are chosen by 5-fold cross-validation on the training items. A tracks the probes on all three tasks (Table~\ref{tab:main}): on prime it lies between the one-hot linear probe (\primeProbeOneHot\%) and the best probe (\primeProbe\%), on div7 both are at chance (\divsevenA\% vs.\ \divsevenProbe\%), and on valid both are near ceiling (\validA\% vs.\ \validProbe\%).

\paragraph{The primality shortcut.} A's answers (recovered from its per-item correctness) coincide with the rule ``the last digit is 1, 3, 7 or 9'' on \primeAagreeLastDigitLo\% and \primeAagreeLastDigitHi\% of test items in seeds 0 and 1, and that rule alone scores \ruleLastDigit\%, within a point of A's accuracy. On this test set A has learned divisibility by 2 and 5, the two divisors visible in the last digit, and nothing else: it calls \primeAprimesCorrect\% of primes prime and gets only \primeAcompCorrect\% of composites right. On the same test set, ``odd'' scores \ruleOdd\%, ``no factor 2 or 3'' \ruleNoTwoThree\% and ``no factor $\le 7$'' \ruleNoSeven\%; A sits between the first two, and a probe with engineered divisibility features does better than A (\primeProbe\%). On the earlier study's independent test set, answer training gave \priorPrimeA\%. The test set contains \primeHardNegN{} odd composites with no prime factor $\le 7$, and A misclassifies every one of them in both seeds (correct: \primeHardNegAcorrect). Such hard negatives are \hardNegFrac\% of all 4-digit composites, so a rule that ignores them costs little on the training distribution. The preregistered test S6 evaluates A on a dedicated hard-negative set of \primeHardN{} items: accuracy \primeHardA\% (B: \primeHardB\%; S6 \verdictSSix).

\paragraph{More answers do not buy the program.} On prime, A reaches \primeAsmall\%, \primeA\% and \primeAlarge\% at $n=\nTrainSmall$, \nTrain{} and \nTrainLarge{}, and \primeAtok\% with the full pool of \primeAtokN{} examples. On div7, A is at chance at every $n$ (\divsevenAsmall\%, \divsevenA\% and \divsevenAlarge\% at $n=\nTrainSmall$, \nTrain{} and \nTrainLarge), and what it learns is a seed-dependent constant answer: seed 0 answers ``no'' to most items (\divsevenAyessZero\% correct on multiples of 7, \divsevenAnosZero\% on non-multiples), seed 1 ``yes'' (\divsevenAyessOne\% and \divsevenAnosOne\%). A is also below the base model on div7 (\divsevenBase\%; McNemar $p=\divsevenABasepsZero$ and $\divsevenABasepsOne$): the base model answers after long free-form responses (its evaluation on prime takes \baseEvalRatio{} times as long as A's), and answer-only training removes them. Part of what a trace adds may therefore be test-time computation; B$'$ and the prompting controls (\S\ref{sec:controls}) separate the two.

\paragraph{The gain is label correctness.} With labels permuted across training items, the earlier study's answer arm sat at \shufMin--\shufMax\% on all seven of its tasks, against up to \priorCorrectMax\% with correct labels (Appendix~\ref{app:tables}). Correct labels are what drive A; on primality, what correct labels buy is the rule.

\subsection{Exact traces teach short programs}
\label{sec:traces}
\paragraph{Divisibility by 7.} B reaches \divsevenB\% against \divsevenA\% for A (discordant items B/A: \divsevenBApsZeroDisc; $p=\divsevenBApsZero$) and \divsevenC\% for C ($p=\divsevenBCpsZero$). The gain holds on both classes, \divsevenByes\% on multiples of 7 and \divsevenBno\% on non-multiples, so it is not a constant answer. At $n=\nTrainSmall$ the same arm was at \divsevenBsmall\%; the effect appears between \nTrainSmall{} and \nTrain{} examples, with one seed at each size. The preregistered replication (S1) gives \divsevenBsOne\% and \divsevenBsTwo\% for seeds 1 and 2 (A seed 2: \divsevenAsTwo\%; McNemar $p=\divsevenBApsOne$ and $\divsevenBApsTwo$; S1 \verdictSOne).

\paragraph{Other divisors (S2, S3).} Holding $k=\divsevenSteps$ and varying $d$, the map predicts that answers suffice for $d=2$, whose label is the last digit, and that traces are needed for $d\in\{3,11,13\}$, which are coprime to 10 (\S\ref{sec:why}). Table~\ref{tab:prereg} gives the results (S2 \verdictSTwo; S3 \verdictSThree). In the earlier answer-only study, $d=3$ gave \priorDivthreeA\% and $d=13$ \priorDivthirteenA\%.

\paragraph{Primality: a long trace hurts.} B reaches \primeB\% against \primeA\% for A ($p=\primeBApsZero$ and $\primeBApsOne$), and B vs.\ C is not significant ($p=\primeBCpsZero$ and $\primeBCpsOne$): on primality the content of the trace gives no detectable benefit over its scrambled tokens. The per-class split locates the failure. A prime's trace must get \primeStepsMin--\primeStepsMax{} consecutive ``$n \bmod p\neq0$'' steps right, and one false zero turns it into a composite; a composite's trace usually ends at its first step. B classifies \primeBprimesCorrect\% of primes and \primeBcompCorrect\% of composites correctly. More traces help the long class: at $n=\nTrainLarge$, B reaches \primeBlarge\% with \primeBlargePrimes\% of primes correct, still below A.

\paragraph{Length (S5).} On 6-digit div7 ($k=\divsevenSixSteps$), the preregistered prediction is B $\approx p^{6}$, with $p=\stepAcc$ measured on 4-digit B generations, i.e.\ \divsevenSixPred. B reaches \divsevenSixB\% (A: \divsevenSixA\%; S5 \verdictSFive). The balanced-test form of the prediction derived in \S\ref{sec:why}, $(1+p^6)/2=\divsevenSixPredBal$, was not preregistered and is reported as a sensitivity analysis.

\paragraph{Validity.} Both formats work: A \validA\%, B \validB\% (\validBsZero\% and \validBsOne\% by seed; B is below A in seed 1, $p=\validBApsOne$), C \validC\%. Because the \nForms{} forms are shared between training and test, this task measures surface learnability, which the probe confirms (Appendix~\ref{app:logic}).

\paragraph{Cost.} B's training corpora contain \tokRatioLow--\tokRatioHigh{} times as many completion tokens as A's (\primeTokB{} vs.\ \tokA{} on prime). Where answers work, they are far cheaper, so the first preregistration's efficiency hypothesis H3 failed on prime and valid; on div7, where answers never leave chance at any $n$ we ran, it holds trivially. In sum, of the first preregistration's endpoints, H1 (B $>$ A) held on div7 only, H2 (B $>$ C) held on div7 and valid but not on prime, H3 failed where answers work, and H4 held at $n=\nTrain$ but not at $n=\nTrainSmall$.

\begin{figure}[t]
\centering
\figorbox{fig_arms.pdf}{0.9\linewidth}{4.6cm}
\caption{Accuracy by arm and task with 95\% Wilson intervals (\nTest{} test items per task). Dashed line: majority class.}
\label{fig:arms}
\end{figure}

\subsection{Controls}
\label{sec:controls}
\paragraph{Content, not tokens.} C sits at \divsevenC\% on div7 and \validC\% on valid, far below B, although it trains on B's tokens with the true answers. What helps on these tasks is the trace's content. On prime, B and C are statistically indistinguishable, and both are below A.

\paragraph{Training signal or test-time computation (B$'$).} In B, the model may simply run the procedure at test time. B$'$ puts the answer first, so the procedure cannot compute it; if B$'$ exceeds A, trace tokens change what the model learns, and if B$'$ matches A, the gain is test-time computation learned from demonstrations. B$'$ on div7: \divsevenBp\%.

\paragraph{Prompting the base model.} A zero-shot CoT prompt gives \primeCot\%, \divsevenCot\% and \validCot\% on prime, div7 and valid, and \nShots{} trace demonstrations in the prompt give \primeFew\%, \divsevenFew\% and \validFew\%. These separate ``the base model runs the procedure when shown it'' from ``fine-tuning on traces taught it''.

\begin{table}[t]
\caption{Preregistered tests (Appendix~\ref{app:prereg}) with results as they came out. Qwen2.5-1.5B-Instruct and 4-digit inputs unless stated; accuracy in \%, mean of two seeds unless stated; McNemar per seed, Holm across tasks.}
\label{tab:prereg}
\centering\footnotesize\setlength{\tabcolsep}{4pt}
\begin{tabular}{@{}llcccp{0.3\linewidth}c@{}}
\toprule
& cell & probe & A & B & criterion & met \\
\midrule
S1 & div7, seeds 1 / 2 & \divsevenProbe & \divsevenAsOne{} / \divsevenAsTwo & \divsevenBsOne{} / \divsevenBsTwo & B $\ge 80$, B$-$A $\ge 20$\pp, $p<0.01$ in each seed & \verdictSOne \\
\addlinespace
S2 & div3 & \divthreeProbe & \divthreeA & \divthreeB & \multirow{3}{=}{$\ge 2$ of the 3 with B$-$A $\ge 15$\pp{} in both seeds} & \multirow{3}{*}{\verdictSTwo} \\
   & div11 & \divelevenProbe & \divelevenA & \divelevenB & & \\
   & div13 & \divthirteenProbe & \divthirteenA & \divthirteenB & & \\
\addlinespace
S3 & div2 & \divtwoProbe & \divtwoA & \divtwoB & A $\ge 95$ & \verdictSThree \\
S4 & div7, 3B model & & \threeBdivsevenA & \threeBdivsevenB & B $>$ A & \verdictSFour \\
S5 & div7, 6-digit & & \divsevenSixA & \divsevenSixB & $|\mathrm{B}-p^6|\le 10$\pp; $p^6=\divsevenSixPred$ & \verdictSFive \\
S6 & prime, hard neg. & & \primeHardA & \primeHardB & A $<60$ and agrees more with ``odd, no factor $\le 7$'' than with the label & \verdictSSix \\
\bottomrule
\end{tabular}
\end{table}

\section{Scale}
\label{sec:scale}
The map locates the answer frontier by the probe, which does not depend on the model, and the trace frontier by $p^k$, which should move with model size. Under the same protocol, Qwen2.5-3B-Instruct gives, as A/B: prime \threeBprimeA/\threeBprimeB, div7 \threeBdivsevenA/\threeBdivsevenB, valid \threeBvalidA/\threeBvalidB{} (S4 \verdictSFour). Where compute allowed, div7 at 0.5B gives \halfBdivsevenA/\halfBdivsevenB{} and at 7B \sevenBdivsevenA/\sevenBdivsevenB. The map predicts that A stays near the probe at every size, while B's accuracy on a fixed task rises with $p$; a larger model should also narrow B's deficit on primality only if its per-step accuracy approaches 1, since $k$ reaches \primeStepsMax.

\section{Why: traces turn divisibility into a table}
\label{sec:why}
\begin{remark}[informal]
\label{rem:automaton}
Let $x_1\cdots x_k$ be the decimal digits of $n$ and $\delta_d(r,a)=(10r+a)\bmod d$. Then $d\mid n$ iff $\delta_d^{*}(0,x_1\cdots x_k)=0$, where $\delta_d^{*}$ applies $\delta_d$ digit by digit: divisibility by $d$ is recognized by a $d$-state automaton. The arm-B trace writes the states $r_1,\dots,r_k$. Under teacher forcing, each supervised state $r_t=\delta_d(r_{t-1},x_t)$ is a function of two tokens already in the context, so $n$ training traces are $kn$ samples of a single table with $10d$ entries, whatever $k$ is (\divsevenStepSamples{} samples of \divsevenTableSize{} entries for $d=7$, $k=\divsevenSteps$, $n=\nTrain$). Answer-only supervision gives $n$ samples of the $k$-fold composition, and for $d$ coprime to 10 its value is nearly independent of every proper subset of the digits: for $d=7$ and $k=\divsevenSteps$, fixing any three digits leaves $\Pr(7\mid n)$ between \condLow{} and \condHigh, against a marginal of \condMarg.
\end{remark}
In the terms of \citet{abbe2024far}, the answer-only target is as global as the input allows, while each trace step is local given the previous state; their inductive scratchpad, the teacher-forced CoT of \citet{kim2025parity} and the sub-task decomposition of \citet{wies2023subtask} are settings in which this kind of decomposition provably makes learning efficient. We do not claim the remark extends those results: they concern transformers trained from scratch on specified distributions, and we fine-tune a pretrained model on \nTrain{} examples. The remark says which of our tasks should benefit, and it makes three predictions.

\emph{(i) Surface.} If $d$ divides a power of 10 ($d=2,5$), the label is a function of the last digits, a probe finds it, and answers suffice (S3). \emph{(ii) Table size.} For $d$ coprime to 10 the label depends on every digit, so answers can succeed only through a rule over all digits (the digit sum for $d=3$), and traces should help with a table of $10d$ entries, so at fixed $n$ the benefit may shrink as $d$ grows (S2). \emph{(iii) Compounding.} If each state is right with probability $p$ independently and a wrong state is uniform on $\mathbb{Z}_d$, the answer is right with probability $p^k+(1-p^k)/d$ on multiples of $d$ and $p^k+(1-p^k)(d-1)/d$ on non-multiples, hence $(1+p^k)/2$ on a balanced test set (Appendix~\ref{app:compound}). The preregistered S5 uses $p^k$, which is the prediction for the positive class up to $1/d$.

\paragraph{Primality through the same lens.} A trial-division step asks the model to reduce a four-digit number modulo a prime in one step: a table indexed by $(n,p)$ with more than $10^5$ entries, almost none of which recur in training, rather than a $10d$-entry table seen many times. A prime's trace has \primeStepsMin--\primeStepsMax{} such steps, and the asymmetric per-class split of \S\ref{sec:traces} is the signature (iii) predicts for a long trace in which one error flips a positive into a negative. A digit-level trace for each divisor would make every step local but multiply $k$ by four; the map predicts that this trade fails unless $p$ is very close to 1. We have not run it. \paragraph{Validity.} The trace is a certificate whose first clause (satisfiable, or contradiction) fixes the answer, so $k$ is effectively 1, and the forms are shared with training; both arms work, as both conditions predict.

\section{Discussion and limitations}
\label{sec:discussion}
\paragraph{What exact labels guarantee.} An exact verifier guarantees that no training label is wrong. It does not guarantee that the labels identify the verifier's function among the functions a model can fit from $n$ examples. On primality, exactly correct labels taught a last-digit rule that calls every composite ending in 1, 3, 7 or 9 prime, among them all \hardNegPop{} of the \compPop{} 4-digit composites without a factor $\le 7$. A shuffled-label control shows that the gain comes from label correctness, but it cannot tell the verifier's function from a surface rule; hard negatives chosen from the probe's errors can.

\paragraph{Supervision design for verifiable-reward pipelines.} Three practices follow from the map. (1) Fit a cheap surface probe on the verified data before training; where it is already accurate, answer supervision will teach its function, and the evaluation should include inputs on which the probe is wrong. (2) Where a trace is used, make each step a function of a few tokens already in context and keep $k$ small, or measure the per-step accuracy and budget $p^k$; a correct but long trace can be worse than no trace. (3) Do not read a null from one training size; our own \nTrainSmall-example null on div7 was the edge of a transition, not the absence of an effect.

\paragraph{No RLVR arm, and what the map predicts for one.} With a binary verifier, the RLVR reward is a function of the final answer only, so the map treats RLVR as answer supervision with on-policy sampling. On a surface-learnable task, a surface rule already earns most of the reward on the training distribution (the last-digit rule A found is right on \ruleLastDigit\% of our prime test items, ``no factor $\le 7$'' on \ruleNoSeven\%), and the map predicts that RLVR finds a shortcut of the same kind. On div7 the reward carries no gradient toward the procedure unless the policy's own samples sometimes execute it; the base model's zero-shot \divsevenBase\% shows that its free-form outputs carry some signal, and \citet{yue2025rl} suggest that RLVR mostly reweights such samples. We did not run GRPO or any RL method. It is the first experiment we would add, evaluated on the hard-negative set.

\paragraph{Limitations.} One model family (Qwen2.5), with 1.5B as the main size; LoRA with one configuration. Binary tasks with balanced test sets of \nTest{} items, so effects under about \ciHalf\pp{} are not resolved, and the per-class and hard-negative figures rest on \nTestClass{} and \primeHardNegN{} items. One trace format per task, written by us; the map predicts how other formats should behave, but we tested one. At the time of writing the div7 result at $n=\nTrain$ rests on one seed, and its replication is preregistered (S1). The first preregistration's conditional arm D (a fluent trace from another instance, close to the corrupted traces of \citet{valmeekam2026beyond}) now qualifies on div7 and has not been run. The shuffled-label control comes from an earlier study with a related protocol (answer-only prompts, 3--4-digit numbers). The probe is one choice, defined after the prior grid; a stronger probe would move the answer frontier. The logic task shares \nForms{} forms between training and test and tests surface learnability, not logical reasoning. Fine-tuning on answers also removes the base model's test-time reasoning, a confound that B$'$ and the prompting controls address.

\subsection*{AI use statement}
We used an AI assistant (Claude, Anthropic) to write and refactor experiment code, to run the experiments and log their outputs under our direction, to compute the statistics, to draft and edit the manuscript, and to check every bibliography entry against arXiv, OpenReview, proceedings and DOI records. The research question, the arms, the preregistered decision rules and the interpretation were reviewed and approved by the authors before the runs they govern. We reviewed all AI-assisted code, text and numbers, verified reported numbers against the released result files, and take responsibility for the content of this paper.

\subsection*{Reproducibility statement}
The supplementary material contains the dataset generators (with fixed seeds), the single-run training and evaluation worker, the queue runner, both preregistration files and the commit that froze the second, and every run's result record, including a per-item correctness bitmap from which each paired test in this paper can be recomputed. Section~\ref{sec:setup} and Appendix~\ref{app:formats} give the prompts, trace formats and hyperparameters. All numbers in the text are macros regenerated from the result files by scripts in the \texttt{paper/} directory, including the surface probe and the derived per-class statistics (\texttt{paper/exp/prior\_stats.py}). Each run takes \runMinLow--\runMinHigh{} minutes on one 8\,GB consumer GPU.

\subsection*{Ethics statement}
The study uses synthetic arithmetic and logic data only and involves no human subjects or personal data.

\bibliography{refs}
\bibliographystyle{iclr2027_conference}

\appendix
\section{Full results}
\label{app:tables}
Table~\ref{tab:runs} lists every run of the prior grid as logged. ``B-ext'' uses reasoning traces sampled from a locally served 9B model and kept only when their final answer matched the verifier (STaR-style rejection sampling); only \primeBornN{} traces passed, so these runs are not comparable to B and are listed for completeness. Train tokens count whitespace tokens in the completions. Table~\ref{tab:shuffled} reports the earlier answer-only study with its shuffled-label control.

\begin{table}[h]
\caption{Every run of the prior grid (Qwen2.5-1.5B-Instruct, \nTest{} test items per task).}
\label{tab:runs}
\centering\small
\IfFileExists{tab_runs.tex}{% generated by paper/exp/make_numbers.py from the run files -- do not edit
\begin{longtable}{lllrrrr}
\caption{Every run, as logged (\nTest{} test items per cell). Cell: the test set, prefixed by the training task when they differ; [cot] and [fewshot4] mark prompting controls on the base model. Train tokens: whitespace tokens in the completions.}\label{tab:runs}\\
\toprule
model & cell & arm & $n$ & seed & train tokens & accuracy (\%) \\
\midrule
\endfirsthead
\toprule
model & cell & arm & $n$ & seed & train tokens & accuracy (\%) \\
\midrule
\endhead
\bottomrule
\endfoot
1.5B & prime & base & 0 & 0 & 0 & 59.6 \\
1.5B & prime & A & 60 & 0 & 120 & 59.6 \\
1.5B & prime & A & 180 & 0 & 360 & 82.1 \\
1.5B & prime & A & 180 & 1 & 360 & 82.5 \\
1.5B & prime & A & 540 & 0 & 1{,}080 & 75.0 \\
1.5B & prime & B & 60 & 0 & 4{,}355 & 57.1 \\
1.5B & prime & B & 180 & 0 & 13{,}370 & 57.9 \\
1.5B & prime & B & 180 & 1 & 13{,}260 & 57.5 \\
1.5B & prime & B & 540 & 0 & 40{,}095 & 68.8 \\
1.5B & prime & C & 180 & 0 & 12{,}561 & 50.0 \\
1.5B & prime & C & 180 & 1 & 13{,}120 & 54.2 \\
1.5B & prime & A-full & 1320 & 0 & 2{,}640 & 80.8 \\
1.5B & prime & B-ext & 9 & 0 & 896 & 62.1 \\
1.5B & prime & B-ext & 9 & 1 & 974 & 60.0 \\
\midrule
1.5B & div7 & base & 0 & 0 & 0 & 62.5 \\
1.5B & div7 & A & 60 & 0 & 120 & 50.0 \\
1.5B & div7 & A & 180 & 0 & 360 & 52.1 \\
1.5B & div7 & A & 180 & 1 & 360 & 50.8 \\
1.5B & div7 & A & 180 & 2 & 360 & 47.9 \\
1.5B & div7 & A & 540 & 0 & 1{,}080 & 50.0 \\
1.5B & div7 & B & 60 & 0 & 3{,}600 & 50.0 \\
1.5B & div7 & B & 180 & 0 & 10{,}800 & 92.5 \\
1.5B & div7 & B & 180 & 1 & 10{,}800 & 90.0 \\
1.5B & div7 & B & 180 & 2 & 10{,}800 & 82.5 \\
1.5B & div7 & C & 180 & 0 & 10{,}784 & 49.6 \\
1.5B & div7 & C & 180 & 1 & 10{,}824 & 50.4 \\
\midrule
1.5B & valid & base & 0 & 0 & 0 & 30.0 \\
1.5B & valid & A & 60 & 0 & 120 & 97.5 \\
1.5B & valid & A & 180 & 0 & 360 & 99.2 \\
1.5B & valid & A & 180 & 1 & 360 & 100.0 \\
1.5B & valid & A & 540 & 0 & 1{,}080 & 100.0 \\
1.5B & valid & B & 60 & 0 & 2{,}343 & 87.5 \\
1.5B & valid & B & 180 & 0 & 7{,}001 & 100.0 \\
1.5B & valid & B & 180 & 1 & 7{,}004 & 80.0 \\
1.5B & valid & B & 540 & 0 & 21{,}084 & 100.0 \\
1.5B & valid & C & 180 & 0 & 7{,}092 & 54.2 \\
1.5B & valid & C & 180 & 1 & 7{,}090 & 50.8 \\
1.5B & valid & A-full & 1600 & 0 & 3{,}200 & 100.0 \\
\midrule
1.5B & div13 & A & 180 & 0 & 360 & 49.6 \\
1.5B & div13 & B & 180 & 0 & 10{,}800 & 54.6 \\
\end{longtable}
}{\fbox{Run table pending (run \texttt{paper/exp/prior\_stats.py}).}}
\end{table}

\begin{table}[h]
\caption{Earlier answer-only study (same model and LoRA configuration; \priorTrainN{} training items, answer-only prompts, 3--4-digit numbers for arithmetic; mean of two seeds). Correct labels vs.\ labels permuted across training items.}
\label{tab:shuffled}
\centering\small
\begin{tabular}{lccc}
\toprule
task & base & correct labels & shuffled labels \\
\midrule
divisible by 3 & \priorDivthreeBase & \priorDivthreeA & \priorDivthreeShuf \\
divisible by 7 & \priorDivsevenBase & \priorDivsevenA & \priorDivsevenShuf \\
divisible by 13 & \priorDivthirteenBase & \priorDivthirteenA & \priorDivthirteenShuf \\
perfect square & \priorSquareBase & \priorSquareA & \priorSquareShuf \\
prime & \priorPrimeBase & \priorPrimeA & \priorPrimeShuf \\
argument validity & \priorValidBase & \priorValidA & \priorValidShuf \\
joint consistency (SAT) & \priorConsistBase & \priorConsistA & \priorConsistShuf \\
\bottomrule
\end{tabular}
\end{table}

\section{The logic task}
\label{app:logic}
Arguments instantiate \nFormsHalf{} valid forms (modus ponens, modus tollens, hypothetical syllogism, disjunctive syllogism, constructive dilemma, contraposition) and \nFormsHalf{} invalid ones (affirming the consequent, denying the antecedent, affirming a disjunct, a broken chain, the converse, and ``or'' to ``and'') with distinct random atom letters, rendered in English (``if $A$ then $B$'', ``not $A$'', ``$A$ or $B$''). The schema proposes an argument and Z3 decides it: the argument is labeled valid iff the premises together with the negated conclusion are unsatisfiable. For a valid argument the trace names the unsatisfiable core; for an invalid one it gives the satisfying assignment as a countermodel. Training and test share the \nForms{} forms and never an instance, so a model can reach high accuracy by recognizing the form, which the atom-renamed $n$-gram probe does (\validProbe\%). The task therefore checks that answers work when the surface determines the label; it is not a test of logical reasoning. Its results: base \validBase\%, A \validA\% (\validAsZero{} and \validAsOne), B \validB\% (\validBsZero{} and \validBsOne), C \validC\% (\validCsZero{} and \validCsOne); at $n=\nTrainSmall$, A \validAsmall\% and B \validBsmall\%; at $n=\nTrainLarge$, both \validAlarge\%.

\section{Trace formats}
\label{app:formats}
All arms share the prompt, which ends ``End your reply with `Answer: Yes' or `Answer: No'.'' Examples of training completions, verbatim:

\begin{small}
\noindent\textbf{div7, B} (prompt: ``Is 1617 divisible by 7?''):
\begin{quote}\ttfamily\raggedright
Process the digits of 1617 keeping a running remainder mod 7. r = (10*0 + 1) mod 7 = 1. r = (10*1 + 6) mod 7 = 2. r = (10*2 + 1) mod 7 = 0. r = (10*0 + 7) mod 7 = 0. Final remainder 0, so 1617 is divisible by 7.\\Answer: Yes
\end{quote}
\noindent\textbf{div7, C} (a scrambled trace from another example; the answer line keeps this example's label):
\begin{quote}\ttfamily\raggedright
running mod r = so = 4067 Process by 7. 7 a = 0. 0) 0. 7 mod (10*5 = r 7 the = = keeping + r is mod mod 4) (10*0 + of 4067 r + 7) 5. = + digits remainder 7. (10*0 = 7 0, divisible Final 6) 4. (10*4 mod remainder\\Answer: Yes
\end{quote}
\noindent\textbf{prime, B}, a composite (``Is 4143 a prime number?''):
\begin{quote}\ttfamily\raggedright
Check prime divisors up to isqrt(4143) = 64. 4143 mod 2 = 1. 4143 mod 3 = 0. 3 divides 4143, so 4143 is composite.\\Answer: No
\end{quote}
\noindent\textbf{prime, B}, a prime (``Is 1493 a prime number?''):
\begin{quote}\ttfamily\raggedright
Check prime divisors up to isqrt(1493) = 38. 1493 mod 2 = 1. 1493 mod 3 = 2. 1493 mod 5 = 3. 1493 mod 7 = 2. 1493 mod 11 = 8. 1493 mod 13 = 11. 1493 mod 17 = 14. 1493 mod 19 = 11. 1493 mod 23 = 21. 1493 mod 29 = 14. 1493 mod 31 = 5. 1493 mod 37 = 13. No prime divisor up to 38, so 1493 is prime.\\Answer: Yes
\end{quote}
\noindent\textbf{valid, B}, invalid (``Premises: (1) if H then D; (2) D. Conclusion: H.''):
\begin{quote}\ttfamily\raggedright
Assume every premise is true and the conclusion false. That is satisfiable: set D=true, H=false. Under this assignment all premises hold while the conclusion is false, a countermodel, so the argument is invalid.\\Answer: No
\end{quote}
\noindent\textbf{valid, B}, valid (``Premises: (1) if S then R; (2) S. Conclusion: R.''):
\begin{quote}\ttfamily\raggedright
Assume every premise is true and the conclusion false, i.e. \textasciitilde(R) holds. Together these force a contradiction: \{S -> R, S, \textasciitilde(R)\} cannot all hold at once. Since assuming the conclusion false is impossible, the conclusion follows.\\Answer: Yes
\end{quote}
\noindent\textbf{A} (every task): \texttt{Answer: Yes} or \texttt{Answer: No}.
\end{small}

\section{Preregistration text}
\label{app:prereg}
\paragraph{First preregistration (prior grid).} Frozen before any run of the prior grid and never edited; results were appended below it. It fixed: arms A, B, C, an answers-only arm at larger scale, arms using model-generated traces, and a conditional wrong-trace arm; Qwen2.5-1.5B-Instruct with the LoRA configuration of \S\ref{sec:setup}; tasks prime, valid and div7, all reported; \nTest{} problem-disjoint test items per task (\nTestClass{} per class), identical across arms, graded by one parser on the final answer. Endpoints: H1, B $>$ A at $n=\nTrain$ per task, paired McNemar, Holm across the three tasks; H2, B $>$ C at $n=\nTrain$; H3, accuracy at matched training tokens; H4, on div7 B reaches about 60\% with A at chance. The full text is in the supplementary material.

\paragraph{Second preregistration (this paper's new runs).} Reproduced from \texttt{PREDICTIONS.md}; Markdown markup converted, text unchanged except one redacted location.
\begin{small}
\begin{quote}
\textbf{Preregistered predictions and decision rules (frozen before any run listed here).} Committed 2026-09-25 \textasciitilde20:35 [location redacted], before any of the runs below were started. The commit hash of this file is the preregistration. Earlier results (results/THINKING\_VS\_DATA.md, frozen 2026-08-04, and its later staged results block) are prior data, reported as such.

\textbf{Provenance of the prior div7 result (stated in the paper).} The 2026-08-04 verdict ``H4 (a procedure trace unlocks div7): NO'' was written from the n = 60 trace run alone (B@60 = 50.0\%). The preregistered main cell, B@180 seed 0, finished afterwards at 92.5\% with the same worker (no code change). The paper reports the premature verdict and the later result.

\textbf{Claim under test.} Exact labels teach shortcuts; exact traces teach short programs. With labels from an exact verifier:
answer-only supervision (arm A) recovers roughly what a simple probe on the answer surface recovers;
trace supervision (arm B: procedure trace, then answer) succeeds when the per-step trace accuracy p, compounded over k steps, stays high (predicted accuracy about p\textasciicircum k), and fails when the procedure is long.

\textbf{Decision rules.}
\textbf{S1 (go/no-go, before any other new run):} div7 B@180 seeds 1 and 2 each reach >= 80\% accuracy, with B - A >= 20 percentage points and McNemar p < 0.01 against A of the same seed. If S1 fails, the paper is not submitted in this form.
\textbf{S2:} at least 2 of \{div3, div11, div13\} (4-digit numbers) give B - A >= 15 pp with the same sign in both seeds.
\textbf{S3:} div2 A >= 95\% (the answer is surface-learnable from the last digit).
\textbf{S4:} the div7 sign (B > A) holds for Qwen2.5-3B-Instruct.
\textbf{S5 (falsifiable length test):} on 6-digit div7, B's accuracy is within 10 pp of p\textasciicircum 6, where p is B's per-step accuracy measured on 4-digit div7 generations.
\textbf{S6 (shortcut):} A on prime agrees more with the rule ``odd and no factor <= 7'' than with the true label on hard negatives (odd composites with no factor <= 7), where A's accuracy drops below 60\%.

Definitions: a ``win'' is a difference of >= 10 pp with both seeds agreeing in sign; ``neither'' is both arms at or below majority class + 8 pp. McNemar tests per seed, Holm correction across tasks within each hypothesis family, alpha = 0.05. Every cell is reported, including those that go against the claim.

\textbf{Arms.} base (zero-shot), A (answers), B (trace then answer), C (length-matched scrambled trace), B' (answer then trace, div7 only), plus eval-only controls: zero-shot chain-of-thought prompt and 4-shot trace demonstrations on the base model. Model: Qwen2.5-1.5B-Instruct (laptop), Qwen2.5-3B-Instruct and others where the RTX 5090 allows. LoRA r = 8, alpha 16, 3 epochs, lr 2e-4, bf16, greedy decoding, n = 180 training examples, 240 problem-disjoint test items.
\end{quote}
\end{small}

\section{The compounding approximation}
\label{app:compound}
Let an item have true final state $r^\star_k$ and suppose each emitted state is correct with probability $p$, independently, and that after the first error the emitted final state is uniform on $\mathbb{Z}_d$. With probability $p^k$ the trace is error-free and the answer is right. Otherwise the answer is right if the emitted final state falls on the correct side of the test $r_k=0$: with probability $1/d$ when $r^\star_k=0$ (a multiple of $d$) and $(d-1)/d$ otherwise. On a balanced test set the accuracy is therefore $p^k+\tfrac12(1-p^k)\,[\tfrac1d+\tfrac{d-1}d]=(1+p^k)/2$. Two assumptions can fail: errors can be correlated across steps (a model that miscomputes one table entry does so whenever the entry recurs), and a wrong state need not be uniform. The positive-class accuracy, $p^k+(1-p^k)/d$, is the cleaner test of compounding, because a model that answers ``no'' by default also scores well on non-multiples: on 6-digit div7 it is \divsevenSixByes\%.

\end{document}
